\emph{Setting.} By Lemma \claimref{mass.de-sitter-scale-frozen-point} and part (i) of Proposition \claimref{synthesis.frozen-phase-spectrum-not-red}, the frozen phase in flat slicing is $g = \eta/(Hx_0)^{2}$, and in the usual sign conventions the metric is $a^{2}(d\eta^{2} - d\vec x^{2})$ with $a = -1/(H\eta)$, $\eta = x_0 < 0$, the late-time boundary being $\eta \to 0$. A massless minimally coupled scalar (weight $0$ in the dictionary $m^{2} = -w(w + 3)H^{2}$ used there) with $\chi = a\phi$ obeys $\chi'' + (k^{2} - a''/a)\chi = 0$ with $a''/a = 2/\eta^{2}$; in $z = k\eta$ this is $\chi_{zz} + (1 - 2/z^{2})\chi = 0$. The Bunch-Davies solution is $u(z) = e^{-I_{4} z}(1 - I_{4}/z)$ (times $(2k)^{-1/2}$), standard in de Sitter field theory (Birrell and Davies, cited from memory); the script verifies that it solves the equation, as an identity in $z$, and that $u\bar u' - u'\bar u = 2I_{4}$.

\emph{(i)} Write $\chi = \alpha u + \beta\bar u$ (units $(2k)^{-1/2} = 1$) and impose $\chi = 1$, $\chi_z = -I_{4}$ at $z_d = k\eta_d = -x$. With $W = 2I_{4}$, $\alpha = (\bar u'(z_d) + I_{4}\bar u(z_d))/W$ and $\beta = -(u'(z_d) + I_{4} u(z_d))/W$. Since $u' = e^{-I_{4} z}(-I_{4} - 1/z + I_{4}/z^{2})$, $u' + I_{4} u = e^{-I_{4} z}\,I_{4}/z^{2}$ and $\bar u' + I_{4}\bar u = e^{I_{4} z}(2I_{4} - 2/z - I_{4}/z^{2})$, so $\alpha = e^{I_{4} z_d}(1 + I_{4}/z_d - 1/(2z_d^{2}))$ and $\beta = -e^{-I_{4} z_d}/(2z_d^{2})$; at $z_d = -x$ these are the displayed $\alpha = e^{-I_{4} x}A$, $\beta = -e^{I_{4} x}B$. Then $|\alpha|^{2} - |\beta|^{2} = |A|^{2} - B^{2} = 1 + 1/(4x^{4}) - 1/(4x^{4}) = 1$. As $\eta \to 0$, $u \to -I_{4}/z$ and $\bar u \to I_{4}/z$ at leading order, so $\chi \to (-I_{4}/z)(\alpha - \beta)$ and $|\phi|^{2} = |\chi|^{2}/a^{2}$ is $|\alpha - \beta|^{2}$ times its Bunch-Davies value, whose dimensionless spectrum is $(H/2\pi)^{2}$. Finally $|\alpha - \beta|^{2} = |A + e^{2I_{4} x}B|^{2} = |A|^{2} + B^{2} + 2B\,\mathrm{Re}(\bar A e^{2I_{4} x})$, and $\mathrm{Re}(\bar Ae^{2I_{4} x}) = (1 - 1/(2x^{2}))\cos 2x - x^{-1}\sin 2x$, which gives $F$. The script stores every expression as $e^{\mp I_{4} z}$ times a Laurent polynomial in $z^{-1}$ with Gaussian-rational coefficients and differentiates exactly; it checks the mode equation, the Wronskian, the coefficients $A$ and $B$, the normalisation, and the closed form of $F$ as an identity in independent $(x, \cos 2x, \sin 2x)$ on rational points of the circle (only $|e^{2I_{4} x}| = 1$ is used).

\emph{(ii)} $x^{4}F(x) = x^{4} + \tfrac12 + (x^{2} - \tfrac12)\cos 2x - x\sin 2x$; inserting the Taylor series of $\cos 2x$ and $\sin 2x$ with their rational coefficients, the coefficients of $x^{0}, \dots, x^{5}$ cancel and those of $x^{6}, x^{8}, x^{10}$ are $\tfrac49, -\tfrac1{15}, \tfrac{8}{1575}$ (computed exactly by the script). So $F = \tfrac49 x^{2}(1 - \tfrac{3}{20}x^{2} + O(x^{4}))$ and $d\ln F/d\ln x \to 2$.

\emph{(iii)} $F = |A + \omega B|^{2}$ with $|\omega| = 1$ lies between $(|A| - B)^{2}$ and $(|A| + B)^{2}$. Using $1 \le |A| \le 1 + 1/(8x^{4})$ (from $\sqrt{1 + t} \le 1 + t/2$), $F(x_1) \le (1 + 1/(8x_1^{4}) + 1/(2x_1^{2}))^{2}$, decreasing in $x_1$, and $F(x_2) \ge (1 - 1/(2x_2^{2}))^{2}$. A constant index $n_s$ over $[x_1, x_2]$, $x_2 = e^{7}x_1$, needs $F(x_1)/F(x_2) = e^{7(1 - n_s)} \ge e^{7 \times 0.0267}$ at the weakest tilt of the two-standard-deviation band ($1 - n_s \ge 0.0351 - 0.0084$), and $e^{y} \ge 1 + y + y^{2}/2$. At $x_1 = 3$, $x_2 \ge 1096 \cdot 3$ ($e^{7} > 1096$ from a partial sum of the exponential series), the upper bound of the ratio is about $1.118$ against the lower bound $1.205$ of the required factor; the script checks this inequality exactly, and that the same bound at $x_1 = 2$ would not suffice, so the cut at $3$ is where the exact argument stops. The window width is the measured multipole range $2 \lesssim \ell \lesssim 2500$, $\ln 1250 \approx 7.13$ (Planck 2018, cited as in Proposition \claimref{synthesis.frozen-phase-spectrum-not-red}); the measured index and its error are those used there.

\emph{(iv)} Every window, exactly. A constant index $n_s$ in the band on $[x_1, x_2]$, $x_2 = e^{7}x_1 > 1096x_1$, needs $F(x_1)/F(x_2) \ge e^{7 \times 0.0267} \ge 1 + y + y^{2}/2 =: R$, $y = 7 \times 0.0267$. Three exact enclosures of $F$ suffice. (a) On $(0, 1]$: $x^{4}F = \sum_n s_nx^{n}$ with $s_0 = \dots = s_5 = 0$, $s_6 = \tfrac49$, and $|s_n| \le 3 \cdot 2^{n}/(n-2)!$ (each of the three contributions to the coefficient of $x^{n}$ in $(x^{2} - \tfrac12)\cos 2x - x\sin 2x$ is at most $2^{n}/(n-2)!$), so $F(x) \le Sx^{2}$ and $F(x) \ge (\tfrac49 - T)x^{2}$ with $S = \sum_{n \ge 6}|s_n| < 0.53$ and $T = \sum_{n \ge 7}|s_n|$, $\tfrac49 - T > 0.37$ (partial sums to $n = 30$ plus the geometric tail bound, ratio at most $1/15$). (b) On $[1, 3]$: write $F = 1 + 1/(2x^{4}) + G/x^{2}$ with $G = (1 - 1/(2x^{2}))\cos 2x - x^{-1}\sin 2x$, $|G'| \le 6$ for $x \ge 1$; on the cells $[x_0, x_0 + \tfrac1{100}]$, with $\cos 2x_0$, $\sin 2x_0$ enclosed by degree-$40$ Taylor sums and the Lagrange remainder, $0.30 < F < 1.16$. (c) On $[3, \infty)$ the envelope bounds of (iii). Then: for $x_1 \ge 1$, $F(x_1) < 1.16$ and $F(x_2) \ge (1 - 1/(2 \cdot 1096^{2}))^{2}$; for $x_1 \le 1$ and $x_2 \ge 3$, $F(x_1) \le S$ and $F(x_2) \ge (1 - \tfrac1{18})^{2}$; for $x_1 \le 1 \le x_2 \le 3$, $x_1 \le 3/1096$ and $F(x_1) \le Sx_1^{2}$ against $F(x_2) > 0.30$; for $x_2 \le 1$, $F(x_1)/F(x_2) \le S/((\tfrac49 - T)1096^{2})$. In every case the ratio is below $R$; the script checks each inequality in exact rationals.

Corroboration, numerical, double precision, not part of the check: on a logarithmic grid of $12000$ starting points $x_1 \in [0.05, 5 \times 10^{4}]$, the mean index $\ln(F(x_1)/F(e^{7}x_1))/7$ is at most $0.0175$ (at $x_1 \approx 2.745$); $F$ has its global maximum $1.1302$ at $x \approx 2.7435$; least-squares fits of $\ln F$ against $\ln x$ on $401$ points per window give a most negative slope of $-0.0047$ (at $x_1 \approx 2.09$) with maximal residual $0.10$. 

\emph{The closing paragraph} reads (ii) to (iv) together: modes outside at the descent are suppressed, modes near the descent scale ring and modes well inside are flat, and no window of the measured width has the measured red drop. A reading, argued and not proved: a state fixed at a single time with no second scale depends on $k$ only through $k|\eta_d|$ and tends to the flat spectrum for $k|\eta_d| \gg 1$, which suggests that a descent producing the tilt would need dilatation breaking lasting while the observed window is set (about seven e-folds of the frozen phase if horizon exit sets each mode); other initial states (adiabatic vacua of higher order) change $A$ and $B$ by terms of higher order in $1/x$ and are not computed here. Transport content: none in the computation, which is standard de Sitter field theory; the setting (frozen phase, massless point, descent at the origin) is the series'. Which de Sitter scale the frozen phase has during the descent is not addressed.
